Sparse hidden signals are a realistic failure mode for text monitoring, but most prior evaluations of linguistic steganography focus on dense generation schemes, local edit rules, or binary cover-versus-stego classification. We study a narrower and practically important question: can current large language models detect and recover messages when the signal is diluted across long passages rather than concentrated in obvious local patterns? We build a controlled benchmark from \wikitext, \imdb, and \agnews with three encoding schemes: acrostics, dense keyword-trigger rules, and diluted keyword-trigger rules with distractor variants. We then evaluate \gptfourone and \gptfive with fixed detection and extraction prompts on 72 synthetic examples, plus a small external sanity check.

The main result is consistent across both models. Exact recovery drops sharply under dilution, while detection degrades more slowly. For \gptfourone, exact recovery falls from 0.75 on dense keyword encoding and 0.67 on acrostics to 0.25 on diluted keyword encoding. For \gptfive, the same ordering appears at 0.50, 0.58, and 0.17. The strongest confirmatory comparison is for \gptfourone, where dense extraction significantly outperforms diluted extraction ($p=0.039$ by Fisher's exact test).

These results show that current LLM-based monitors can often notice suspicious structure without reliably recovering the payload once the signal is sparse, dispersed, and mixed with ordinary text. That detection-extraction gap matters for misuse analysis and for benchmarking model-based steganalysis systems.
